% Options for packages loaded elsewhere
\PassOptionsToPackage{unicode}{hyperref}
\PassOptionsToPackage{hyphens}{url}
\PassOptionsToPackage{dvipsnames,svgnames,x11names}{xcolor}
\documentclass[
  10pt,
  fontset=fandol]{ctexart}
\usepackage{xcolor}
\usepackage[margin=16mm]{geometry}
\usepackage{amsmath,amssymb}
\setcounter{secnumdepth}{-\maxdimen} % remove section numbering
\usepackage{iftex}
\ifPDFTeX
  \usepackage[T1]{fontenc}
  \usepackage[utf8]{inputenc}
  \usepackage{textcomp} % provide euro and other symbols
\else % if luatex or xetex
  \usepackage{unicode-math} % this also loads fontspec
  \defaultfontfeatures{Scale=MatchLowercase}
  \defaultfontfeatures[\rmfamily]{Ligatures=TeX,Scale=1}
\fi
\usepackage{lmodern}
\ifPDFTeX\else
  % xetex/luatex font selection
\fi
% Use upquote if available, for straight quotes in verbatim environments
\IfFileExists{upquote.sty}{\usepackage{upquote}}{}
\IfFileExists{microtype.sty}{% use microtype if available
  \usepackage[]{microtype}
  \UseMicrotypeSet[protrusion]{basicmath} % disable protrusion for tt fonts
}{}
\makeatletter
\@ifundefined{KOMAClassName}{% if non-KOMA class
  \IfFileExists{parskip.sty}{%
    \usepackage{parskip}
  }{% else
    \setlength{\parindent}{0pt}
    \setlength{\parskip}{6pt plus 2pt minus 1pt}}
}{% if KOMA class
  \KOMAoptions{parskip=half}}
\makeatother
\setlength{\emergencystretch}{3em} % prevent overfull lines
\providecommand{\tightlist}{%
  \setlength{\itemsep}{0pt}\setlength{\parskip}{0pt}}
\usepackage{amsmath,amssymb,mathtools,needspace}
\xeCJKDeclareCharClass{CJK}{"2032,"0400->"04FF,"2160->"216B,"2460->"2473,"25A0,"25C7,"25CB,"25CF}
\IfFontExistsTF{Arial Unicode MS}{\xeCJKsetup{AutoFallBack=true}\setCJKfallbackfamilyfont{\CJKrmdefault}{Arial Unicode MS}}{}
\setcounter{secnumdepth}{0}
\setlength{\emergencystretch}{2em}
\allowdisplaybreaks
\usepackage{bookmark}
\IfFileExists{xurl.sty}{\usepackage{xurl}}{} % add URL line breaks if available
\urlstyle{same}
\hypersetup{
  pdftitle={什么是并行计算},
  colorlinks=true,
  linkcolor={Maroon},
  filecolor={Maroon},
  citecolor={Blue},
  urlcolor={Blue},
  pdfcreator={LaTeX via pandoc}}

\title{什么是并行计算}
\author{}
\date{}

\begin{document}
\maketitle

\subsection{11.1　什么是并行计算}\label{ux4ec0ux4e48ux662fux5e76ux884cux8ba1ux7b97}

\subsubsection{11.1.1　一则寓言故事}\label{ux4e00ux5219ux5bd3ux8a00ux6545ux4e8b}

20 世纪 80 年代初国产银河巨型机问世，国内掀起一股并行算法热。究竟什么是并行计算呢？这里先讲一个生动的寓言故事。

相传很久很久以前，有一个年轻的国王名叫川行，是个数学天才。川行爱上了邻国聪颖美丽且爱好数学的公主邱比郑南。

川行差人前往邻国求婚。公主答应了这桩婚事，但提出了一项先决条件，她要亲自考核一下川行的数学才能。公主的考题是，针对一个 15 位数求出它的真因子。

接到试题之后，川行立即忙碌起来，一个数接一个数地试算。川行有数学才能，算得很快，然而由于 15 位数的真因子可能是个 8 位数，找出全部真因子要花费上亿次

\href{../../source-images/pdf-277.jpeg}{原页图像}

整数除法，总的计算量大得惊人。

川行感到很为难。这是一道``大数分解''的数学难题，如何才能尽快地找出它的答案呢？

川行有个足智多谋的宰相名叫孔幻士。孔幻士提出了一个计谋：将全国老百姓按军、师、团、营、连、排、班、兵 8 个等级编号，每 10 个兵组成 1 班，10 个班为 1 排，10 个排为 1 连\ldots\ldots10 个师为 1 军，10 个军全归川行统帅。这样，在编的每个老百姓都有一个 8 位十进制的编号。完成这种编制以后，通知全国老百姓用自己的编号去除公主给出的 15 位数，能除尽的立即上报，给予重奖。这样很快找出了所有的真因子，而川行则依靠全国老百姓的帮助赢得了公主的爱情。

这则寓言浅显易懂，但意味深长。公主``邱比郑南''是``求比证难''的谐音。大数分解问题的可解性不言而喻，但具体求解却很困难。国王``川行''是``串行''的谐音。串行计算的效率很低，往往不能承担大规模的计算工程。宰相``孔幻士''则是``空换时''的谐音，其含义是，并行计算的设计思想是用扩大空间、增加处理机台数为代价来换取计算时间的节省。``空换时''是并行计算的基本策略。

这则故事所涉及的大数分解问题有重大的学术价值，求解这类问题的计算量随着``大数''的增大而急剧增长。譬如，计算一个 155 位数的真因子，如果用串行算法进行计算，即使用每秒亿次的巨型机去承担，也得要连续工作上万年。当然这是没有实际意义的。1990 年 6 月 20 日美国报道了一则消息：贝尔实验室用 1 000 台处理机并行计算，仅仅花费了几个月的时间，就成功地找出一个 155 位数的 3 个真因子，它们分别是 7 位数、49 位数和 99 位数。这是科学计算的一项重大突破。当年我国《科技日报》评价这项成就为``1990 年世界十大科技成就之一''。

\subsubsection{11.1.2　同步并行算法的设计策略}\label{ux540cux6b65ux5e76ux884cux7b97ux6cd5ux7684ux8bbeux8ba1ux7b56ux7565}

采取并行处理方式运行的计算机系统称作并行机系统，简称并行机。

并行机出现于 20 世纪的 70 年代初，至今仅有近 50 年的历史。1972 年，美国研制成功阵列机 Illiac IV，此后于 1976 年又进一步研制出向量机 Cray-1。并行机的更新换代和商业化开发强有力地推动了并行计算的蓬勃发展。

并行机的体系结构各不相同，但大致可分为两类：一类是单指令流多数据流 SIMD（single instruction stream, multiple data stream）型，如阵列机、向量机；另一类是多指令流多数据流 MIMD（multiple instruction stream, multiple data stream）型，称作多处理机。

针对 SIMD 与 MIMD 两类并行机系统，并行算法大致分为同步与异步两类。本章仅研究同步并行算法。

所谓同步性，是指不同处理机在同一时刻针对不同数据执行同一种操作。同步并行计算的典型例子是向量计算。

\href{../../source-images/pdf-278.jpeg}{原页图像}

同步并行计算的基本策略是``分而治之''。所谓分而治之，就是将所考察的计算问题分裂成若干较小的子问题，并将这些子问题映射到多台处理机上去各自完成，然后再将分散的结果拼装成所求的解。

\textbf{值得指出的是，在同步并行算法设计时，``分而治之''的设计原则往往被误解为整体上的先分后治，而将``分''与``治''两个环节截然分开。这种算法设计技术就是所谓的倍增技术。}

\textbf{其实，在并行计算过程中，``分''与``治''是矛盾的两个方面，它们既是对立的，又是统一的。基于这种理解我们推荐了同步并行算法设计的二分技术。}

需要强调的是，为避免局限于具体的机器特征而束缚了并行算法的研究，人们提出了理想计算机的概念。``理想化''的假设包括：任何时刻可以使用任意多台处理机，任何时刻有任意多个主存单元可供使用，处理机同主存间的数据通信时间可以忽略不计。

\end{document}
